\section*{الفصل \ref{ch:g12:seq} --- المتتاليات}

\begin{solution}{exo:g12:seq:1}
لتكن $P(n)$ العبارة $\sum_{k=1}^{n} k^2 = \frac{n(n+1)(2n+1)}{6}$.
\emph{الحالة الابتدائية:} من أجل $n = 0$ يساوي الطرفان $0$ (مجموع خالٍ).
\emph{خطوة التراجع:} لنفترض $P(n)$. عندئذٍ
\[
\sum_{k=1}^{n+1} k^2 = \frac{n(n+1)(2n+1)}{6} + (n+1)^2
= \frac{(n+1)\bigl(n(2n+1) + 6(n+1)\bigr)}{6}
= \frac{(n+1)(2n^2 + 7n + 6)}{6}.
\]
وبما أن $2n^2 + 7n + 6 = (n+2)(2n+3)$، فهذا هو
$\frac{(n+1)(n+2)(2(n+1)+1)}{6}$، أي $P(n+1)$. وبالتراجع، تصح $P(n)$
من أجل كل $n$.
\end{solution}

\begin{solution}{exo:g12:seq:2}
$a_{n+1} - a_n = \frac{n+2}{n+1} - \frac{n+1}{n}
= \frac{n(n+2) - (n+1)^2}{n(n+1)} = \frac{-1}{n(n+1)} < 0$: إذن $(a_n)$
\omterm{def:g10:functions:variations}{متناقصة} تمامًا.

وحدود $(b_n)$ موجبة و
$\frac{b_{n+1}}{b_n} = \frac{2^{n+1}}{n+1}\cdot\frac{n}{2^n} =
\frac{2n}{n+1} \geq 1 \iff 2n \geq n+1 \iff n \geq 1$: إذن $(b_n)$ \omterm{def:g12:seq:monotonic}{متزايدة}
(تمامًا من أجل $n \geq 2$).

$c_{n+1} - c_n = (n+1)^2 - 10(n+1) - n^2 + 10n = 2n - 9$، وهو سالب
من أجل $n \leq 4$ وموجب من أجل $n \geq 5$: فتتناقص $(c_n)$ حتى
$c_5 = -25$، وهي قيمتها الصغرى، ثم تتزايد. فهي ليست \omterm{def:g12:seq:monotonic}{رتيبة}.
\end{solution}

\begin{solution}{exo:g12:seq:3}
أخرج الحدود المهيمنة عوامل مشتركة:
\[
u_n = \frac{n^2(3 - 1/n + 1/n^2)}{n^2(2 + 5/n^2)} \xrightarrow[n\to+\infty]{} \frac{3}{2}.
\]
اضرب في المرافق:
\[
v_n = \frac{(n+1) - n}{\sqrt{n+1} + \sqrt{n}} = \frac{1}{\sqrt{n+1}+\sqrt{n}}
\xrightarrow[n\to+\infty]{} 0.
\]
اقسم البسط والمقام على $3^n$:
\[
w_n = \frac{(2/3)^n - 1}{1 + (1/3)^n} \xrightarrow[n\to+\infty]{} \frac{0-1}{1+0} = -1,
\]
باستعمال $\lim q^n = 0$ من أجل $\abs{q} < 1$.
\end{solution}

\begin{solution}{exo:g12:seq:4}
$u_n = 5 + 3n \to +\infty$ و $v_n = 8 \cdot (1/2)^n = 2^{3-n} \to 0$.
والمجموعان هما
\[
\sum_{k=0}^{n} u_k = (n+1)\,\frac{5 + (5+3n)}{2} = \frac{(n+1)(10+3n)}{2}
\xrightarrow[n\to+\infty]{} +\infty,
\]
\[
\sum_{k=0}^{n} v_k = 8\,\frac{1 - (1/2)^{n+1}}{1 - 1/2}
= 16\left(1 - \left(\tfrac{1}{2}\right)^{n+1}\right)
\xrightarrow[n\to+\infty]{} 16 .
\]
\end{solution}

\begin{solution}{exo:g12:seq:5}
بما أن $-1 \leq \cos n \leq 1$، فإن
\[
\frac{n-1}{n+1} \leq \frac{n + \cos n}{n+1} \leq 1,
\]
و $\frac{n-1}{n+1} \to 1$، إذن النهاية هي $1$ بمبرهنة الحصر.

وأما النهاية الثانية، فاكتب
\[
0 \leq \frac{n!}{n^n}
= \frac{1}{n}\cdot\frac{2}{n}\cdots\frac{n}{n}
\leq \frac{1}{n},
\]
لأن كل عامل $\frac{k}{n}$ حيث $2 \leq k \leq n$ لا يتجاوز $1$.
وبما أن $\frac1n \to 0$، تعطي مبرهنة الحصر أن $\frac{n!}{n^n} \to 0$. (والحصر
الهندسي المقترح يفلح أيضًا: فكل عامل حيث $k \leq n/2$ لا يتجاوز
$\frac12$، مما يعطي الحاصرة الأقوى $(1/2)^{\floor{n/2}}$.)
\end{solution}

\begin{solution}{exo:g12:seq:6}
\emph{1.} $u_0 = 0 \in \intcc{0}{2}$. وإذا كان $0 \leq u_n \leq 2$، فإن
$2 \leq u_n + 2 \leq 4$، إذن $\sqrt{2} \leq u_{n+1} \leq 2$؛ وخاصة
$0 \leq u_{n+1} \leq 2$. وبالتراجع تصح الخاصية من أجل كل $n$.

\emph{2.} $u_{n+1} - u_n = \sqrt{u_n + 2} - u_n$. ومن أجل $x \in \intcc{0}{2}$، يكون
$\sqrt{x+2} \geq x \iff x + 2 \geq x^2 \iff (2-x)(x+1) \geq 0$، وهذا صحيح.
ومنه فإن $(u_n)$ \omterm{def:g12:seq:monotonic}{متزايدة}.

\emph{3.} وبما أنها \omterm{def:g12:seq:monotonic}{متزايدة} \omterm{def:g12:seq:bounded}{ومحدودة من الأعلى} بالعدد $2$، \omterm{def:g12:seq:limit}{تتقارب} $(u_n)$ إلى
$\ell \in \intcc{0}{2}$ ما. وبالانتقال إلى النهاية في $u_{n+1} = \sqrt{u_n + 2}$
(فالتطبيق $x \mapsto \sqrt{x+2}$ متصل) نجد $\ell = \sqrt{\ell + 2}$،
إذن $\ell^2 - \ell - 2 = 0$، أي $\ell \in \{-1, 2\}$. وبما أن $\ell \geq 0$،
فإن $\lim u_n = 2$.
\end{solution}

\begin{solution}{exo:g12:seq:7}
\emph{1.} بين جرعتين، يُطرح $40\%$ من الدواء، فتصير
الكمية $u_n$ هي $0.6\,u_n$؛ ثم تضيف الجرعة التالية $1$ وحدة:
$u_{n+1} = 0.6\,u_n + 1$.

\emph{2.} $v_{n+1} = u_{n+1} - 2.5 = 0.6\,u_n + 1 - 2.5 = 0.6(u_n - 2.5)
= 0.6\,v_n$: إذن $(v_n)$ \omterm{def:g12:seq:arith-geom}{هندسية} أساسها الهندسي $0.6$ وحدها الأول
$v_0 = 1 - 2.5 = -1.5$. ومنه $v_n = -1.5 \times 0.6^n$ و
\[
u_n = 2.5 - 1.5 \times 0.6^n .
\]

\emph{3.} وبما أن $0.6^n \to 0$، فإن $u_n \to 2.5$: أي أن كمية الدواء تستقر
عند $2.5$ وحدة.
\end{solution}

\begin{solution}{exo:g12:seq:8}
\emph{1.} $u_0 = 3 > 1$. وإذا كان $u_n > 1$، فإن $u_n + 2 > 0$ و
\[
u_{n+1} - 1 = \frac{4u_n - 1 - u_n - 2}{u_n + 2} = \frac{3(u_n - 1)}{u_n + 2} > 0 .
\]
وبالتراجع، يكون $u_n > 1$ من أجل كل $n$ (وخاصة $u_n + 2 \neq 0$، إذن
\omterm{def:g12:seq:sequence}{المتتالية} معرَّفة جيدًا).

\emph{2.} باستعمال المتطابقة أعلاه،
\[
v_{n+1} = \frac{1}{u_{n+1} - 1} = \frac{u_n + 2}{3(u_n - 1)}
= \frac{(u_n - 1) + 3}{3(u_n - 1)} = \frac{1}{3} + v_n .
\]
إذن $(v_n)$ \omterm{def:g12:seq:arith-geom}{حسابية} أساسها $\frac13$ و
$v_0 = \frac{1}{u_0 - 1} = \frac12$.

\emph{3.} $v_n = \frac12 + \frac{n}{3}$، ومنه
$u_n = 1 + \frac{1}{v_n} = 1 + \frac{6}{3 + 2n}$. وبما أن $v_n \to +\infty$،
فإن $u_n \to 1$.
\end{solution}

\begin{solution}{exo:g12:seq:9}
\emph{1.} المقدار $H_{2n} - H_n = \sum_{k=n+1}^{2n} \frac{1}{k}$ مجموع $n$
حدًا، كل منها لا يقل عن $\frac{1}{2n}$؛ ومنه
$H_{2n} - H_n \geq n \cdot \frac{1}{2n} = \frac12$.

\emph{2.} بالتراجع على $k$: $H_{2^0} = H_1 = 1 \geq 1$. وإذا كان
$H_{2^k} \geq 1 + \frac{k}{2}$، فإن تطبيق النقطة 1 مع $n = 2^k$ يعطي
\[
H_{2^{k+1}} \geq H_{2^k} + \frac12 \geq 1 + \frac{k+1}{2}.
\]
\omterm{def:g12:seq:sequence}{والمتتالية} $(H_n)$ \omterm{def:g12:seq:monotonic}{متزايدة} (فكل خطوة تضيف $\frac{1}{n+1} > 0$) و
\omterm{def:g12:seq:sequence}{المتتالية} الجزئية $H_{2^k}$ غير \omterm{def:g12:seq:bounded}{محدودة}، إذن $(H_n)$ غير \omterm{def:g12:seq:bounded}{محدودة من الأعلى}.
وبما أنها \omterm{def:g12:seq:monotonic}{متزايدة} وغير \omterm{def:g12:seq:bounded}{محدودة}، فهي تؤول إلى $+\infty$
(\cref{thm:g12:seq:monotone}).
\end{solution}

\begin{solution}{exo:g12:seq:10}
\emph{1.} \omterm{def:g12:seq:sequence}{المتتالية} $d_n = b_n - a_n$ تحقق
$d_{n+1} - d_n = (b_{n+1} - b_n) - (a_{n+1} - a_n) \leq 0$، إذن $(d_n)$
\omterm{def:g10:functions:variations}{متناقصة}؛ وبما أن $d_n \to 0$، نجد $d_n \geq 0$ من أجل كل $n$ (\omterm{def:g12:seq:sequence}{فمتتالية}
\omterm{def:g10:functions:variations}{متناقصة} فيها حد سالب تبقى دونه إلى الأبد، مما يمنع
النهاية $0$). ومنه $a_n \leq b_n$.

\emph{2.} من $a_n \leq b_n \leq b_0$، تكون \omterm{def:g12:seq:sequence}{المتتالية} \omterm{def:g12:seq:monotonic}{المتزايدة} $(a_n)$
\omterm{def:g12:seq:bounded}{محدودة من الأعلى}، \omterm{def:g12:seq:limit}{فتتقارب} إلى $\ell$ ما. وبالمثل \omterm{def:g12:seq:limit}{تتقارب} $(b_n)$، \omterm{def:g10:functions:variations}{المتناقصة}
\omterm{def:g12:seq:bounded}{والمحدودة من الأسفل} بالعدد $a_0$، إلى $\ell'$ ما. ثم
$\ell' - \ell = \lim (b_n - a_n) = 0$، إذن $\ell = \ell'$.

\emph{3.} \omterm{def:g12:seq:sequence}{المتتالية} $(a_n)$ \omterm{def:g12:seq:monotonic}{متزايدة} (تمامًا) لأن
$a_{n+1} - a_n = \frac{1}{(n+1)!} > 0$. وأما $(b_n)$، فإن
\[
b_{n+1} - b_n = \frac{1}{(n+1)!} + \frac{1}{(n+1)(n+1)!} - \frac{1}{n\,n!}
= \frac{n(n+1) + n - (n+1)^2}{n(n+1)(n+1)!}
= \frac{-1}{n(n+1)(n+1)!} < 0 ,
\]
إذن $(b_n)$ \omterm{def:g10:functions:variations}{متناقصة}. وأخيرًا $b_n - a_n = \frac{1}{n\,n!} \to 0$. فالمتتاليتان
متجاورتان، ومنه تتقاربان إلى نهاية مشتركة.
\end{solution}

\begin{solution}{pb:g12:seq:1}
\textbf{1.} صحيحة من أجل $n = 1$ ($1 = 1^2$). وإذا كان
$1 + 3 + \dots + (2n - 1) = n^2$، فبإضافة العدد الفردي
التالي: $n^2 + (2n + 1) = (n + 1)^2$: وهي الوراثة. وبالتراجع،
تصح من أجل كل $n \geq 1$.

\textbf{2.} $2^0 = 1 > 0$. وإذا كان $2^n > n$، فإن
$2^{n+1} = 2 \cdot 2^n > 2n \geq n + 1$ من أجل $n \geq 1$ (والحالة
$n = 0$ تُتحقق مباشرة): وهي الوراثة، وقد تمّ المطلوب.

\textbf{3.} $n = 0$: $1 \geq 1$. وإذا كان
$(1 + x)^n \geq 1 + nx$، فاضرب في $1 + x \geq 1 > 0$:
$(1 + x)^{n+1} \geq (1 + nx)(1 + x) = 1 + (n + 1)x + nx^2
\geq 1 + (n + 1)x$.

\textbf{4.} الخطوة من $n = 1$ إلى $n = 2$: ففي حالة كريتين،
تكون ``الكريات $n$ الأولى'' و``الكريات $n$ الأخيرة'' كريتين
مفردتين \emph{منفصلتين} --- فلا كرية مشتركة تجسر بين
المجموعتين، ولا شيء يفرض توافق لونيهما. وحجة
الوراثة تقتضي ضمنًا أن تتقاطع المجموعتان،
وهذا لا يصح إلا من $n \geq 2$؛ ومع الحالة الابتدائية $n = 1$ لا تنطلق
السلسلة أبدًا.

\textbf{5.} $4^0 - 1 = 0 = 3 \times 0$. وإذا كان $4^n - 1 = 3k$،
فإن $4^{n+1} - 1 = 4(4^n - 1) + 3 = 3(4k + 1)$: وهي الوراثة.

\textbf{6.} $x_1 = \frac32$، و $x_2 = \frac{17}{12}$،
و $x_3 = \frac{577}{408}$.

\textbf{7.} $x_{n+1}^2 - 2 = \frac{(x_n^2 + 2)^2 - 8x_n^2}
{4x_n^2} = \frac{(x_n^2 - 2)^2}{4 x_n^2}$: مربع على
مقدار موجب، ومنه فهو $\geq 0$، وهو $> 0$ كلما كان $x_n^2 \neq 2$.
والتراجع: $x_0 = 2 > 0$ مع $x_0^2 = 4 > 2$؛ وإذا كان $x_n > 0$
و $x_n^2 > 2$، فإن $x_{n+1}$ (وهو \omterm{def:g10:stats:mean}{متوسط} موجبين)
موجب و $x_{n+1}^2 - 2 > 0$.

\textbf{8.} $x_{n+1} - x_n = \frac{2 - x_n^2}{2x_n} < 0$ حسب
السؤال 7: إذن \omterm{def:g10:functions:variations}{متناقصة} تمامًا.

\textbf{9.} \omterm{def:g10:functions:variations}{متناقصة} \omterm{def:g12:seq:bounded}{ومحدودة من الأسفل} (بالعدد $1$، لأن
$x_n^2 > 2 > 1$ و $x_n > 0$): فبمبرهنة التقارب
الرتيب، \omterm{def:g12:seq:limit}{تتقارب} $(x_n)$ إلى $L \geq 1$ ما.

\textbf{10.} تحترم النهايات الجبر: فمن
$x_{n+1} = \frac12\left(x_n + \frac{2}{x_n}\right)$ و
$x_n \to L \geq 1 > 0$:
$L = \frac12\left(L + \frac2L\right)$، إذن $L^2 = 2$، وبما أن $L$
موجب، فإن $L = \sqrt2$. والحكم: التقارب مبرهَن عليه،
والنهاية محدَّدة --- فهيرون بريء مع مرتبة الشرف.

\textbf{11.} $x_{n+1} - \sqrt2 = \frac{x_n^2 - 2\sqrt2\,x_n +
2}{2x_n} = \frac{(x_n - \sqrt2)^2}{2x_n}$: أي بالضبط
$e_{n+1} = \frac{e_n^2}{2x_n}$، و $x_n > \sqrt2$ يعطي
$e_{n+1} \leq \frac{e_n^2}{2\sqrt2}$. فالخطأ مربَّع: أي أن كل خطوة
تضاعف عدد الأرقام العشرية الصحيحة، كما لوحظ منذ
الكتاب السابق.

\textbf{12.} $e_0 \approx 0.5858$، و $e_1 \approx 0.0858$،
و $e_2 \approx 0.00245$، و $e_3 \approx 2.1 \times 10^{-6}$.
والنسب $\frac{e_1}{e_0^2} \approx 0.25 = \frac{1}{2x_0}$؛
$\frac{e_2}{e_1^2} \approx 0.333 = \frac{1}{2x_1}$؛
$\frac{e_3}{e_2^2} \approx 0.353 \approx \frac{1}{2x_2}$: أي المبرهنة
في العمل.

\textbf{13.} $H_{2^{18}} \geq 1 + 9 = 10$: أي نحو
$260\,000$ حد ($2^{18} = 262\,144$) لمجرد تجاوز $10$ ---
فهو تباعد بخطوة السلحفاة (وأما $H_n > 100$ فسيحتاج إلى حدود
أكثر من ذرات أي مكتبة).

\textbf{14.} $S_n = 2 - \frac{1}{2^n}$ (مجموع هندسي)، و
$\frac{1}{2^n} \to 0$ (\cref{thm:g12:seq:geometric})، إذن
$S_n \to 2$: فلوح الشوكولاتة المقضوم بلا انتهاء يؤول إلى
اللوح كله دون أن يبلغه أبدًا --- والآن بلغة النهايات
الرسمية.

\textbf{15.} من أجل $k \geq 2$:
$\frac{1}{k^2} \leq \frac{1}{k(k-1)} = \frac{1}{k-1} -
\frac1k$، إذن
$S_n \leq 1 + \left(1 - \frac1n\right) < 2$: فهي \omterm{def:g12:seq:monotonic}{متزايدة} و
\omterm{def:g12:seq:bounded}{محدودة من الأعلى}، ومنه متقاربة (بالتقارب الرتيب). وقد سمّى أويلر
النهاية فيما بعد: $\frac{\pi^2}{6}$.

\textbf{16.} أيلولة الحدود إلى $0$ \emph{لازمة} لكي
تستقر المجاميع لكنها لا تقرّر شيئًا: فالحدود التوافقية
$\frac1n \to 0$ ومع ذلك تنفجر المجاميع؛ والحدود
$\frac{1}{n^2} \to 0$ والمجاميع \omterm{def:g12:seq:limit}{تتقارب}. فالسؤال كله هو
\emph{بأي سرعة} تموت الحدود --- وتلك نظرية المتسلسلات،
المبنية في الكتب الجامعية.

\textbf{17.} $a_1 = \sqrt2 \approx 1.41421$ و $b_1 = 1.5$؛
و $a_2 \approx 1.45648$، و $b_2 \approx 1.45711$: فتكراران
اثنان يتفقان أصلًا إلى ثلاثة أرقام عشرية --- سرعة مذهلة.

\textbf{18.} $b_{n+1} - a_{n+1} = \frac{a_n + b_n}{2} -
\sqrt{a_n b_n} = \frac{(\sqrt{b_n} - \sqrt{a_n})^2}{2} \geq
0$: فالمتوسطان يبقيان مرتبين. و $(a_n)$ \omterm{def:g12:seq:monotonic}{متزايدة}:
$a_{n+1} = \sqrt{a_n b_n} \geq \sqrt{a_n \cdot a_n} = a_n$؛
و $(b_n)$ \omterm{def:g10:functions:variations}{متناقصة} بتناظر.

\textbf{19.} $\frac{b_{n+1} - a_{n+1}}{b_n - a_n} =
\frac{(\sqrt{b_n} - \sqrt{a_n})^2}
{2(\sqrt{b_n} - \sqrt{a_n})(\sqrt{b_n} + \sqrt{a_n})}
= \frac{\sqrt{b_n} - \sqrt{a_n}}{2(\sqrt{b_n} + \sqrt{a_n})}
\leq \frac12$: فالفجوة تنتصف على الأقل، إذن $b_n - a_n \to 0$؛
ومع السؤال 18، تكون المتتاليتان متجاورتين ولهما نهاية مشتركة
$M(1, 2)$.

\textbf{20.} يعطي التكرار الثالث
$a_3 \approx b_3 \approx 1.456791$: أي $M(1, 2) \approx
1.456791$ في ثلاث لفات من المدور (فالفجوة
\emph{تتربّع} \omterm{def:g10:numbers:approx}{تقريبًا}، مثل فجوة هيرون). \omterm{def:g10:coordgeom:system}{وترتيب} سرعات
متتاليات المسألة، من الأبطأ إلى الأسرع: المجاميع التوافقية (تباعد
جليدي)، ثم المجاميع \omterm{def:g12:seq:arith-geom}{الهندسية} (الخطأ ينتصف عند كل خطوة)،
ثم هيرون \omterm{def:g11:stat:mean}{والمتوسط الحسابي} الهندسي (الخطأ يتربّع عند كل خطوة) --- وسرعة
هذا الأخير الخارقة هي التي أنبأت غاوس بأنه أصاب عرقًا جديدًا
من التحليل.
\end{solution}
